\lectures{12,13}
\title{P, NP, SAT, Poly-Time Reducibility}

\begin{document}

\subsection{Time Complexity}

\textbf{Definition.}
Let $M$ be a deterministic Turing machine decider.
The \textbf{time complexity} of $M$ is the function
$f \colon \mathbb{N} \to \mathbb{N}$ defined so that
$f(n)$ is the maximum number of steps that $M$ takes
on any input of length $n$.

\textbf{Definition.}
For any function $t \colon \mathbb{N} \to \mathbb{R}^+$,
we may define its \textbf{time complexity class} $\mathrm{TIME}(t(n))$
to be the set of all languages
that are decidable by an $O(t(n))$-time Turing machine.

\begin{problem}
    Consider the language $A := \{\str{0}^k\str{1}^k \mid k \in \mathbb{N}_0\}$.
    Show that $A \in \mathrm{TIME}(n \log n)$.
\end{problem}

\begin{solution}
    The following Turing machine, on input $w$,
    decides $A$ in $O(n \log n)$ time.
    \begin{enumerate}
        \item Perform a single scan to check that
        $w \in \{\str{0}^a \str{1}^b \mid a, b \in \mathbb{N}_0\}$.
        \item Repeat the following while there exist
        both \str{0}s and \str{1}s on the tape.
        \begin{enumerate}
            \item[(a)] Perform a single scan
            to check whether the total number of \str{0}s and \str{1}s
            is odd.  If it is, \textsc{Reject}.
            \item[(b)] Perform a single scan, and during this scan,
            cross off every other \str{0} and
            cross off every other \str{1}.
        \end{enumerate}
        \item Perform a single scan to check whether
        the tape is now \emph{completely} empty.
        If it is, \textsc{Accept};
        otherwise, \textsc{Reject}.
    \end{enumerate}
    This algorithm performs, at worst, an $O(n)$-time scan
    $O(\log n)$ times, so its time complexity is $O(n \log n)$, as promised.
    ~ $\blacksquare$
\end{solution}

The bound $A \in \mathrm{TIME}(n \log n)$
cannot be improved---in fact,
any language decidable in $o(n \log n)$ time must be regular!

However, this all relies heavily on
taking our model of computation to be the single-tape Turing machine.

\begin{problem}
    Show that $A$ is decidable in $O(n)$ time
    by a two-tape Turing machine.
\end{problem}

\begin{solution}
    Simply cut and paste all of the \str{0}s
    from the first tape to the second tape.
    Then scan the two tapes in unison
    and check that they contain the same number of characters.
\end{solution}

\subsection{Complexity Relationships Among Models}

Fortunately, the time complexity of multitape TMs
is not so different from their single-tape TM counterparts.

\begin{theorem}{Multitape Time Complexity}
    Every $t(n) \geq n$ time multitape TM
    has an equivalent $O(t(n)^2)$ time single-tape TM.
\end{theorem}

\begin{proof}
    We apply the same mapping from
    multitape TMs $M$ to single-tape TMs $S$
    used in \href{/18.404/lectures/6/#multi-tape-turing-machines}{Lecture 6}.

    $S$ simulates each step of $M$ with the following substeps:
    \begin{enumerate}
        \item Read the entire active portion
        $\str{T}_1\str{\#}\str{T}_2\str{\#}\dots\str{\#}\str{T}_k$
        of the tape.
        \item For each subtape $\str{T}_i$,
        simulate $M$ on $\str{T}_i$ and
        write out the result, performing right shifts as needed.
    \end{enumerate}
    The point is that the length of
    $\str{T}_1\str{\#}\str{T}_2\str{\#}\dots\str{\#}\str{T}_k$
    is at most $O(t(n))$,
    so $S$ simulates each step of $M$ in $O(t(n))$ time.
    Since $M$ must perform $O(t(n))$ such steps,
    the time complexity of $S$ is $O(t(n)^2)$.
    ~ $\blacksquare$
\end{proof}

We can perform a similar comparison between
deterministic and nondeterministic TMs.

\textbf{Definition.}
Given a nondeterministic TM $N$ that is a decider,
the \textbf{running time} of $N$ is the function
$f \colon \mathbb{N} \to \mathbb{N}$ defined so that
$f(n)$ is the maximum number of steps that $N$ takes
\emph{on any branch} on any input of length $n$.

\begin{theorem}{Nondeterministic Time Complexity}
    Every $t(n) \geq n$ time nondeterministic (single-tape) TM
    has an equivalent $2^{O(t(n))}$ time deterministic (single-tape) TM.
\end{theorem}

\begin{proof}
    We apply the same mapping from
    nondeterministic TMs $N$ to deterministic TMs $M$
    used in \href{/18.404/lectures/6/#nondeterministic-turing-machines}{Lecture 6}.

    Note that $M$ performs BFS on
    the computation tree of $N$.
    This tree has depth at most $t(n)$,
    so it has at most $2^{O(t(n))}$ nodes.
    Simulating each layer takes $2^{O(t(n))}$ time,
    so the time complexity is $t(n) \cdot 2^{O(t(n))} = 2^{O(t(n))}$,
    as promised. ~ $\blacksquare$
\end{proof}

\subsection{Polynomial Time (P)}

\textbf{Definition.}
Let $\mathrm{P}$ denote the class of languages that are decidable
in polynomial time on a deterministic single-tape TM; that is,
\[ \mathrm{P} := \bigcup_{k = 0}^{\infty} \mathrm{TIME}(n^k). \]

One can show languages are in $\mathrm{P}$ via
typical 6.1220 strategies.  For example,
\begin{itemize}
    \item The language $\mathrm{PATH} := \{
        \langle G, s, t \rangle \mid
        G \text{ is a directed graph with a path from } s \text{ to } t
    \}$ is in $\mathrm{P}$.
    \item The language $\mathrm{RELPRIME} := \{
        \langle x, y \rangle \mid \gcd(x, y) = 1
    \}$ is in $\mathrm{P}$.
\end{itemize}

Perhaps less obviously, one can show the following.

\begin{theorem}{$A_{\mathrm{CFG}} \in \mathrm{P}$}
    The language $A_{\mathrm{CFG}} := \{
        \langle G, w \rangle \mid
        G \text{ is a CFG that generates } w
    \}$ is in $\mathrm{P}$.
\end{theorem}

\begin{proof}
    Recall our strategy for deciding $A_{\mathrm{CFG}}$
    from \href{/18.404/lectures/7/#algorithms-on-grammars-and-turing-machines-a-cfg-e-cfg-eq-cfg-a-tm}{Lecture 7};
    in summary, the strategy was:
    \begin{enumerate}
        \item Convert $G$ to Chomsky Normal Form.
        \item Note that any derivation of the input $w$
        must apply a rule exactly $2|w| - 1$ times;
        search through all such derivations.
    \end{enumerate}
    Step 1 is doable in polynomial time, but Step 2 is not!
    There are exponentially many such derivations.

    The trick, however, is to apply \emph{dynamic programming};
    if $w = \str{w}_1\str{w}_2\dots\str{w}_n$, then
    we memoize the following values:
    \[
    T[i, j] := \{A \mid A \text{ is a variable in } G
    \text{ such that } A \text{ generates }
    \str{w}_i\str{w}_{i + 1}\dots \str{w}_j\}.
    \]
    The key observation is that
    $A$ generates $\str{w}_i\str{w}_{i + 1}\dots \str{w}_j$
    if and only if \dots
    \begin{itemize}
        \item (``Inductive Case'': $i < j$)
        \dots there is a rule $A \to BC$ and some $k \in [i, j)$
        such that $B \in T[i, k]$ and $C \in T[k + 1, j]$.
        \item (``Base Case'': $i = j$)
        \dots there exists the rule $A \to \str{w}_i$ in $G$.
    \end{itemize}
    Therefore, $T[1, n]$ can be computed in $O(n^3|G|)$ time
    via dynamic programming, as desired. ~ $\blacksquare$
\end{proof}

Then the following immediately follows.

\begin{theorem}{$\mathrm{CFL} \subseteq \mathrm{P}$}
    Every context-free language is a member of $\mathrm{P}$.
\end{theorem}

\begin{proof}
    Suppose we wish to show the CFL with CFG $G$ is in $\mathrm{P}$.
    Then on input $w$, we need a poly-time decider
    for whether $w \in L(G)$.
    Such a decider can just
    run the poly-time decider for $A_{\mathrm{CFG}}$
    on input $\langle G, w \rangle$.
    ~ $\blacksquare$
\end{proof}

\subsection{Polynomial Verifiability (NP)}

\textbf{Definition.}
A \textbf{verifier} for a language $A$ is an algorithm $V$
such that:
\[
    A = \{ w \mid V \text{ accepts } \langle w, c \rangle
    \text{ for some string } c \}.
\]
A \textbf{polynomial time verifier} is a verifier $V$
that runs in polynomial time in the length of $w$.

\begin{remark}
    Since $V$ runs in polynomial time in $|w|$,
    it can only read polynomially many symbols of $c$.
    So we can assume that the length $|c|$ of the certificate
    is polynomial in $|w|$.
\end{remark}

\textbf{Definition.}
Let $\mathrm{NP}$ denote the class of languages
that have polynomial time verifiers.

One can prove the following,
which provides an equivalent alternative definition for $\mathrm{NP}$.

\begin{theorem}{NP Iff NTM}
    A language is in $\mathrm{NP}$
    if and only if it is decided
    by some nondeterministic polynomial time TM.
\end{theorem}

\begin{proof}
    Just have the NTM simulate the verifier by guessing the certificate.
    Conversely, a verifier can take an accepting branch of the NTM as the certificate.
    ~ $\blacksquare$
\end{proof}

\textbf{Definition.}
Denote $\mathrm{NTIME}(t(n)) := \{
    L \mid L \text{ is a language decided by an }
    O(t(n)) \text{ time NTM}
\}$.  Then we have:
\[
    \mathrm{NP} = \bigcup_{k = 0}^{\infty} \mathrm{NTIME}(n^k).
\]
Via typical \href{/6.1220/lectures/14/}{6.1220} arguments,
one can demonstrate the following:
\begin{itemize}
    \item The language $\mathrm{CLIQUE} := \{
        \langle G, k \rangle \mid
        G \text{ is an undirected graph with a } k\text{-clique}
    \}$ is in $\mathrm{NP}$.
    \item The language $\mathrm{SUBSET}\text{-}\mathrm{SUM} := \{
        \langle S, t \rangle \mid
        S \text{ is a multiset with a sub-multiset whose sum is } t
    \}$ is in $\mathrm{NP}$.
\end{itemize}

\subsection{Poly-Time Reduction}

Recall how we defined reductions in computability,
back in \href{/18.404/lectures/9/#mapping-reducibility-eq-tm-eq-tm}{Lecture 9}.

\textbf{Definition.}
We say that a language $A$ is \textbf{polynomial-time reducible}
to a language $B$ (denoted $A \leq_p B$) if
there is a polynomial-time computable function
$f \colon \Sigma^* \to \Sigma^*$ such that
$w \in A$ if and only if $f(w) \in B$.

\begin{theorem}{Poly-Time Reducibility and P}
    If $A \leq_p B$ and $B \in \mathrm{P}$,
    then $A \in \mathrm{P}$.
\end{theorem}

\begin{proof}
    A poly-time decider for $A$, on input $w$,
    can compute $f(w)$ in poly-time,
    then run the poly-time decider for $B$ on $f(w)$.
    ~ $\blacksquare$
\end{proof}

\end{document}